\section{Introduction}
\label{sec:introduction}

Quantum machine learning (QML) promises expressivity advantages for certain data regimes,
yet reproducible benchmarks with rigorous statistical testing remain scarce.
We introduce \texttt{quantun-ia}, a modular research lab spanning 20 experiments
from 2D synthetic classification to UCI tabular, MNIST-subset, and sequence tasks.

\paragraph{Contribution.}
We provide (1) a holdout-fair comparison protocol that splits data before preprocessing,
logs multi-seed metrics with bootstrap confidence intervals, and applies Holm-corrected
Wilcoxon tests; (2) an open benchmark suite with parameter-matched classical baselines
and documented negative results; and (3) a reproducible publication pipeline that
exports figures, LaTeX tables, and Zenodo-archived release bundles from append-only logs.

Our headline narrative follows \emph{Option~C --- honest QML benchmark}:
exp\_011--014 on real data, hybrid and robustness studies (exp\_002, exp\_017),
and openly reported failures (exp\_003, exp\_009).
Innovation tracks (adaptive LR, NAS, fusion) are deferred to follow-up papers.

Supporting contributions:
\begin{itemize}
  \item Hypothesis-first experiment workflow with append-only logging.
  \item Parameter-matched classical baselines for fair quantum comparisons.
  \item Open publication pipeline: figures, LaTeX tables, and Zenodo-archived releases.
\end{itemize}

Negative results (curriculum, self-play, entanglement ablations) are documented
in \texttt{docs/negative\_results.md} and strengthen the credibility of positive findings.
